\clearpage
\VerdierPhysicalPageAnchor{0123}%
\begingroup
\small
\setlength{\emergencystretch}{1.5em}%
\setlength{\abovedisplayskip}{4pt plus 2pt minus 2pt}%
\setlength{\belowdisplayskip}{4pt plus 2pt minus 2pt}%

\VerdierSemanticAnchor{ch2-1-2-14}
\noindent\textbf{Remarque 1.2.14.}\quad
Les propriétés étudiées au numéro 1.2 n'ont fait à aucun moment intervenir
l'axiome TRIV. Elles sont donc valables pour toutes les $\mathbb Z$-catégories
additives strictes munies d'un ensemble de triangles vérifiant les axiomes
TRI, TRII et TRIII.

\VerdierSemanticAnchor{ch2-1-3}
\medskip
\noindent\textbf{1.3. Exemples de catégories triangulées.}

\VerdierSemanticAnchor{ch2-1-3-1}
\medskip
\noindent\textbf{1.3.1.}\quad
Soient $C$ une catégorie additive, $K(C)$ la catégorie des complexes de $C$
à homotopie près (chap. I, 2.5.7). La catégorie $K(C)$ est une
$\mathbb Z$-catégorie stricte munie d'une famille de triangles, les triangles
distingués (chap. I, 3.3.1), qui possède les propriétés TRI, TRII et TRIII
(chap. I, 3.3).

\VerdierSemanticAnchor{ch2-1-3-2}
\medskip
\noindent\textbf{Proposition 1.3.2.}\quad
\textit{La catégorie $K(C)$, munie de la famille des triangles distingués,
est une catégorie triangulée.}

\VerdierSemanticAnchor{ch2-1-3-2-triv-reduction}
Il reste à démontrer que la famille des triangles distingués de $K(C)$
possède la propriété TRIV. Soient donc :

\VerdierSemanticAnchor{ch2-1-3-2-commutative-triangle}
\[
  \begin{tikzcd}[column sep=4.2em,row sep=2.6em]
    & X_2 \arrow[dr,"u_1"] & \\
    X_1 \arrow[ur,"u_3"] \arrow[rr,"u_2"'] && X_3
  \end{tikzcd}
\]
\noindent un diagramme commutatif et :

\VerdierSemanticAnchor{ch2-1-3-2-input-triangles}
\[
  \begin{aligned}
    X_1&\xrightarrow{u_3}X_2\xrightarrow{v_3}Z_3
      \xrightarrow{w_3}X_1[1],\\
    X_2&\xrightarrow{u_1}X_3\xrightarrow{v_1}Z_1
      \xrightarrow{w_1}X_2[1],\\
    X_1&\xrightarrow{u_2}X_3\xrightarrow{v_2}Z_2
      \xrightarrow{w_2}X_1[1]
  \end{aligned}
\]
\noindent trois triangles distingués. D'après (chap. I, 3.3.5), il existe
trois triangles distingués :

\VerdierSemanticAnchor{ch2-1-3-2-alternative-triangles}
\[
  \resizebox{0.98\linewidth}{!}{$
    (X_1,X_2,Z'_3,u_3,v'_3,w'_3),\qquad
    (X_2,X_3,Z'_1,u_1,v'_1,w'_1),\qquad
    (X_1,X_3,Z'_2,u_2,v'_2,w'_2)
  $}
\]
\noindent et deux morphismes :

\VerdierSemanticAnchor{ch2-1-3-2-morphisms-boundary}
\[
  m'_1\colon Z'_3\longrightarrow Z'_2,\qquad
  m'_3\colon Z'_2\longrightarrow Z'_1,
\]
\endgroup
